\textbf{Phần 1. Trắc nghiệm lựa chọn}
\setcounterpageref{ex}{0}

\begin{ex} %[5D7V39]
Một mảnh đất hình chữ nhật có chiều dài hơn chiều rộng $18\text{ m}$. Chiều rộng bằng $\dfrac{4}{7}$ chiều dài. Chu vi của mảnh đất đó là
\choice
{$66\text{ m}$}
{\True $132\text{ m}$}
{$120\text{ m}$}
{$144\text{ m}$}
\loigiai{
Hiệu số phần bằng nhau là:
$7 - 4 = 3\text{ (phần)}$.\\
Chiều rộng mảnh đất là:
$18 : 3 \times 4 = 24\text{ (m)}$.\\
Chiều dài mảnh đất là:
$24 + 18 = 42\text{ (m)}$.\\
Chu vi mảnh đất là:
$(42 + 24) \times 2 = 132\text{ (m)}$.\\
\textbf{Đáp án của câu hỏi là B}
}
\end{ex}

\begin{ex} %[5D7V39]
Hiện nay mẹ hơn con $28$ tuổi. Sau $3$ năm nữa, tuổi mẹ gấp $3$ lần tuổi con. Tuổi mẹ hiện nay là
\choice
{$36\text{ tuổi}$}
{\True $39\text{ tuổi}$}
{$42\text{ tuổi}$}
{$45\text{ tuổi}$}
\loigiai{
Hiệu số tuổi của mẹ và con luôn không đổi theo thời gian và bằng $28$ tuổi.\\
Sau $3$ năm nữa, tuổi của con là:
$28 : (3 - 1) \times 1 = 14\text{ (tuổi)}$.\\
Tuổi con hiện nay là:
$14 - 3 = 11\text{ (tuổi)}$.\\
Tuổi mẹ hiện nay là:
$11 + 28 = 39\text{ (tuổi)}$.\\
\textbf{Đáp án của câu hỏi là B}
}
\end{ex}

\begin{ex} %[5D7V39]
Hiệu của hai số là $135$. Nếu viết thêm chữ số $0$ vào bên phải số bé thì được số lớn. Số bé là
\choice
{$12$}
{\True $15$}
{$18$}
{$20$}
\loigiai{
Viết thêm chữ số $0$ vào bên phải số bé thì được số lớn nên số lớn gấp $10$ lần số bé.\\
Hiệu số phần bằng nhau là:
$10 - 1 = 9\text{ (phần)}$.\\
Số bé là:
$135 : 9 = 15$.\\
\textbf{Đáp án của câu hỏi là B}
}
\end{ex}

\begin{ex} %[5D7V39]
Hai thùng chứa nước mắm, thùng thứ nhất nhiều hơn thùng thứ hai $24\text{ l}$. Nếu bớt ở mỗi thùng đi $6\text{ l}$ thì số nước mắm ở thùng thứ hai bằng $\dfrac{2}{5}$ thùng thứ nhất. Số nước mắm ở thùng thứ nhất lúc đầu là
\choice
{$40\text{ l}$}
{\True $46\text{ l}$}
{$50\text{ l}$}
{$52\text{ l}$}
\loigiai{
Khi bớt ở cả hai thùng cùng một lượng như nhau thì hiệu số lít nước mắm không đổi và bằng $24\text{ l}$.\\
Hiệu số phần bằng nhau là:
$5 - 2 = 3\text{ (phần)}$.\\
Số nước mắm ở thùng thứ nhất sau khi bớt là:
$24 : 3 \times 5 = 40\text{ (l)}$.\\
Lúc đầu thùng thứ nhất chứa:
$40 + 6 = 46\text{ (l)}$.\\
\textbf{Đáp án của câu hỏi là B}
}
\end{ex}

\begin{ex} %[5D7V39]
Một thửa ruộng hình chữ nhật có chiều dài hơn chiều rộng $25\text{ m}$. Tỉ số giữa chiều rộng và chiều dài là $\dfrac{2}{3}$. Diện tích của thửa ruộng đó là
\choice
{$2\,500\text{ m}^2$}
{$3\,000\text{ m}^2$}
{\True $3\,750\text{ m}^2$}
{$4\,500\text{ m}^2$}
\loigiai{
Hiệu số phần bằng nhau là:
$3 - 2 = 1\text{ (phần)}$.\\
Chiều rộng thửa ruộng là:
$25 \times 2 = 50\text{ (m)}$.\\
Chiều dài thửa ruộng là:
$50 + 25 = 75\text{ (m)}$.\\
Diện tích thửa ruộng là:
$75 \times 50 = 3\,750\text{ (m}^2\text{)}$.\\
\textbf{Đáp án của câu hỏi là C}
}
\end{ex}

\begin{ex} %[5D7V39]
Lớp 5A có số học sinh nam nhiều hơn số học sinh nữ là $6$ bạn. Biết số học sinh nữ bằng $80\%$ số học sinh nam. Số học sinh của cả lớp 5A là
\choice
{$48\text{ bạn}$}
{\True $54\text{ bạn}$}
{$60\text{ bạn}$}
{$45\text{ bạn}$}
\loigiai{
Đổi: $80\% = \dfrac{4}{5}$.\\
Hiệu số phần bằng nhau là:
$5 - 4 = 1\text{ (phần)}$.\\
Số học sinh nữ là:
$6 \times 4 = 24\text{ (bạn)}$.\\
Số học sinh nam là:
$24 + 6 = 30\text{ (bạn)}$.\\
Cả lớp có tất cả:
$24 + 30 = 54\text{ (học sinh)}$.\\
\textbf{Đáp án của câu hỏi là B}
}
\end{ex}

\begin{ex} %[5D7V39]
Hiệu của hai số thập phân là $16{,}2$. Biết rằng nếu dời dấu phẩy của số bé sang bên phải một hàng thì được số lớn. Số lớn là
\choice
{$1{,}8$}
{\True $18$}
{$180$}
{$162$}
\loigiai{
Dời dấu phẩy của số bé sang phải một hàng thì được số lớn, nghĩa là số lớn gấp $10$ lần số bé.\\
Hiệu số phần bằng nhau là:
$10 - 1 = 9\text{ (phần)}$.\\
Số bé là:
$16{,}2 : 9 = 1{,}8$.\\
Số lớn là:
$1{,}8 \times 10 = 18$.\\
\textbf{Đáp án của câu hỏi là B}
}
\end{ex}

\begin{ex} %[5D7V39]
Hai bao gạo có hiệu khối lượng là $36\text{ kg}$. Nếu cùng thêm vào mỗi bao $4\text{ kg}$ thì bao thứ hai nặng gấp 3 lần bao thứ nhất. Số ki-lô-gam gạo ở bao thứ nhất lúc đầu là
\choice
{$10\text{ kg}$}
{\True $14\text{ kg}$}
{$18\text{ kg}$}
{$22\text{ kg}$}
\loigiai{
Khi cùng thêm vào mỗi bao $4\text{ kg}$ thì hiệu khối lượng không đổi và bằng $36\text{ kg}$.\\
Khối lượng bao thứ nhất sau khi thêm là:
$36 : (3 - 1) \times 1 = 18\text{ (kg)}$.\\
Khối lượng bao thứ nhất lúc đầu là:
$18 - 4 = 14\text{ (kg)}$.\\
\textbf{Đáp án của câu hỏi là B}
}
\end{ex}

\begin{ex} %[5D7V39]
Một đoàn tàu chở khách có số ghế ngồi hạng phổ thông nhiều hơn hạng thương gia là $120$ ghế. Biết số ghế thương gia bằng $\dfrac{1}{4}$ số ghế phổ thông. Tổng số ghế của cả đoàn tàu là
\choice
{$160\text{ ghế}$}
{$180\text{ ghế}$}
{\True $200\text{ ghế}$}
{$240\text{ ghế}$}
\loigiai{
Hiệu số phần bằng nhau là:
$4 - 1 = 3\text{ (phần)}$.\\
Số ghế hạng thương gia là:
$120 : 3 \times 1 = 40\text{ (ghế)}$.\\
Số ghế hạng phổ thông là:
$40 \times 4 = 160\text{ (ghế)}$.\\
Tổng số ghế của đoàn tàu là:
$40 + 160 = 200\text{ (ghế)}$.\\
\textbf{Đáp án của câu hỏi là C}
}
\end{ex}

\begin{ex} %[5D7V39]
Hiệu của hai số là $84$. Nếu bớt ở số lớn đi $4$ đơn vị và thêm vào số bé $4$ đơn vị thì số bé bằng $\dfrac{3}{5}$ số lớn. Số lớn ban đầu là
\choice
{$180$}
{$190$}
{\True $194$}
{$200$}
\loigiai{
Khi bớt ở số lớn $4$ đơn vị và thêm vào số bé $4$ đơn vị thì hiệu mới giảm đi:
$4 + 4 = 8\text{ (đơn vị)}$.\\
Hiệu mới là:
$84 - 8 = 76$.\\
Số lớn mới là:
$76 : (5 - 3) \times 5 = 190$.\\
Số lớn ban đầu là:
$190 + 4 = 194$.\\
\textbf{Đáp án của câu hỏi là C}
}
\end{ex}

\begin{ex} %[5D7V39]
Một tủ sách có hai ngăn. Ngăn dưới nhiều hơn ngăn trên $40$ cuốn sách. Nếu chuyển $10$ cuốn từ ngăn trên xuống ngăn dưới thì số sách ngăn trên bằng $\dfrac{1}{3}$ số sách ngăn dưới. Số sách ngăn trên ban đầu là
\choice
{$30\text{ cuốn}$}
{\True $40\text{ cuốn}$}
{$50\text{ cuốn}$}
{$60\text{ cuốn}$}
\loigiai{
Sau khi chuyển, ngăn dưới nhiều hơn ngăn trên số cuốn là:
$40 + 10 \times 2 = 60\text{ (cuốn)}$.\\
Số sách ngăn trên sau khi chuyển là:
$60 : (3 - 1) \times 1 = 30\text{ (cuốn)}$.\\
Số sách ngăn trên ban đầu là:
$30 + 10 = 40\text{ (cuốn)}$.\\
\textbf{Đáp án của câu hỏi là B}
}
\end{ex}

\begin{ex} %[5D7V39]
Hiệu của hai số bằng $45$. Tỉ số của hai số là $\dfrac{4}{9}$. Tổng của hai số đó là
\choice
{$90$}
{$108$}
{\True $117$}
{$135$}
\loigiai{
Hiệu số phần bằng nhau là:
$9 - 4 = 5\text{ (phần)}$.\\
Tổng số phần bằng nhau là:
$9 + 4 = 13\text{ (phần)}$.\\
Tổng của hai số là:
$45 : 5 \times 13 = 117$.\\
\textbf{Đáp án của câu hỏi là C}
}
\end{ex}

\textbf{Phần 2. Trắc nghiệm đúng sai}
\setcounterpageref{ex}{0}

\begin{ex} %[5D7V39]
Cho hai số có hiệu bằng $60$. Biết số bé bằng $\dfrac{2}{5}$ số lớn.
\choiceTF
{\True Hiệu số phần bằng nhau là $3$ phần}
{\True Số bé là $40$}
{\True Số lớn là $100$}
{Tổng của hai số đó bằng $150$}
\loigiai{
\begin{itemchoice}
\item Đúng. Hiệu số phần: $5 - 2 = 3\text{ (phần)}$.
\item Đúng. Số bé: $60 : 3 \times 2 = 40$.
\item Đúng. Số lớn: $40 + 60 = 100$.
\item Sai. Tổng của hai số: $40 + 100 = 140 
eq 150$.
\end{itemchoice}
}
\end{ex}

\begin{ex} %[5D7V39]
Một sân trường hình chữ nhật có chiều dài hơn chiều rộng $24\text{ m}$. Tỉ số giữa chiều rộng và chiều dài là $\dfrac{3}{5}$.
\choiceTF
{\True Chiều rộng của sân trường là $36\text{ m}$}
{\True Chiều dài của sân trường là $60\text{ m}$}
{\True Chu vi của sân trường là $192\text{ m}$}
{Diện tích của sân trường là $2\,000\text{ m}^2$}
\loigiai{
\begin{itemchoice}
\item Đúng. Chiều rộng: $24 : (5 - 3) \times 3 = 36\text{ (m)}$.
\item Đúng. Chiều dài: $36 + 24 = 60\text{ (m)}$.
\item Đúng. Chu vi: $(60 + 36) \times 2 = 192\text{ (m)}$.
\item Sai. Diện tích: $60 \times 36 = 2\,160\text{ (m}^2\text{)} 
eq 2\,000\text{ m}^2$.
\end{itemchoice}
}
\end{ex}

\textbf{Phần 3. Trắc nghiệm trả lời ngắn}
\setcounterpageref{ex}{0}

\begin{ex} %[5D7V39]
Hiệu của hai số là $72$. Tỉ số của hai số đó là $\dfrac{3}{7}$. Tìm số bé.
\shortans{54}
\loigiai{
Hiệu số phần bằng nhau là:
$7 - 3 = 4\text{ (phần)}$.\\
Số bé là:
$72 : 4 \times 3 = 54$.
}
\end{ex}

\begin{ex} %[5D7V39]
Bố hơn con $30$ tuổi. Tuổi bố gấp $4$ lần tuổi con. Hỏi con bao nhiêu tuổi?
\shortans{10}
\loigiai{
Hiệu số phần bằng nhau là:
$4 - 1 = 3\text{ (phần)}$.\\
Tuổi con là:
$30 : 3 = 10\text{ (tuổi)}$.
}
\end{ex}

\begin{ex} %[5D7V39]
Một cuộn dây đỏ dài hơn cuộn dây xanh là $45\text{ m}$. Sau khi cắt bớt ở mỗi cuộn $5\text{ m}$ thì cuộn đỏ gấp $4$ lần cuộn xanh. Hỏi lúc đầu cuộn dây đỏ dài bao nhiêu mét?
\shortans{65}
\loigiai{
Sau khi cắt cùng $5\text{ m}$, hiệu chiều dài hai cuộn không đổi và bằng $45\text{ m}$.\\
Chiều dài cuộn đỏ sau khi cắt là:
$45 : (4 - 1) \times 4 = 60\text{ (m)}$.\\
Lúc đầu cuộn dây đỏ dài:
$60 + 5 = 65\text{ (m)}$.
}
\end{ex}

\begin{ex} %[5D7V39]
Một mảnh vườn hình chữ nhật có chiều dài hơn chiều rộng $15\text{ m}$. Chiều rộng bằng $\dfrac{2}{5}$ chiều dài. Tính diện tích của mảnh vườn đó theo đơn vị mét vuông.
\shortans{250}
\loigiai{
Chiều rộng mảnh vườn là:
$15 : (5 - 2) \times 2 = 10\text{ (m)}$.\\
Chiều dài mảnh vườn là:
$10 + 15 = 25\text{ (m)}$.\\
Diện tích mảnh vườn là:
$25 \times 10 = 250\text{ (m}^2\text{)}$.
}
\end{ex}

\begin{ex} %[5D7V39]
Hai kho thóc A và B, kho A nhiều hơn kho B là $50$ tấn thóc. Biết rằng số thóc ở kho B bằng $75\%$ số thóc ở kho A. Hỏi kho A chứa bao nhiêu tấn thóc?
\shortans{200}
\loigiai{
Đổi: $75\% = \dfrac{3}{4}$.\\
Hiệu số phần bằng nhau là:
$4 - 3 = 1\text{ (phần)}$.\\
Số thóc ở kho A là:
$50 \times 4 = 200\text{ (tấn)}$.
}
\end{ex}

\begin{ex} %[5D7V39]
Hiệu của hai số là $108$. Nếu tăng số lớn thêm $12$ đơn vị thì số lớn gấp $4$ lần số bé. Tìm số lớn ban đầu.
\shortans{148}
\loigiai{
Hiệu mới sau khi tăng số lớn là:
$108 + 12 = 120$.\\
Số bé là:
$120 : (4 - 1) = 40$.\\
Số lớn ban đầu là:
$40 + 108 = 148$.
}
\end{ex}
